% ============================================================
% FIG 10 — Evidence Fusion Flow (Dempster-Shafer / Yager)
% Visual explanation of how each evidence source is converted
% to a mass function and fused via Yager's rule.
% Based on Section IV-D and the Appendix B worked example.
% ============================================================
% PREAMBLE:
%   \usepackage{tikz}
%   \usepackage{xcolor}
%   \usetikzlibrary{arrows.meta, positioning, shapes.geometric, calc}
%
% PLACEMENT: full width
%   \begin{figure*}[t]
%     \centering
%     % ============================================================
% FIG 10 — Evidence Fusion Flow (Dempster-Shafer / Yager)
% Visual explanation of how each evidence source is converted
% to a mass function and fused via Yager's rule.
% Based on Section IV-D and the Appendix B worked example.
% ============================================================
% PREAMBLE:
%   \usepackage{tikz}
%   \usepackage{xcolor}
%   \usetikzlibrary{arrows.meta, positioning, shapes.geometric, calc}
%
% PLACEMENT: full width
%   \begin{figure*}[t]
%     \centering
%     % ============================================================
% FIG 10 — Evidence Fusion Flow (Dempster-Shafer / Yager)
% Visual explanation of how each evidence source is converted
% to a mass function and fused via Yager's rule.
% Based on Section IV-D and the Appendix B worked example.
% ============================================================
% PREAMBLE:
%   \usepackage{tikz}
%   \usepackage{xcolor}
%   \usetikzlibrary{arrows.meta, positioning, shapes.geometric, calc}
%
% PLACEMENT: full width
%   \begin{figure*}[t]
%     \centering
%     % ============================================================
% FIG 10 — Evidence Fusion Flow (Dempster-Shafer / Yager)
% Visual explanation of how each evidence source is converted
% to a mass function and fused via Yager's rule.
% Based on Section IV-D and the Appendix B worked example.
% ============================================================
% PREAMBLE:
%   \usepackage{tikz}
%   \usepackage{xcolor}
%   \usetikzlibrary{arrows.meta, positioning, shapes.geometric, calc}
%
% PLACEMENT: full width
%   \begin{figure*}[t]
%     \centering
%     \input{fig10_evidence_fusion.tex}
%     \caption{Dempster--Shafer evidence fusion pipeline. Each evidence
%     source is converted to a basic belief assignment over
%     $\Theta=\{T,\bar{T}\}$ (Eq.~4). Sources are fused sequentially
%     via Yager's conflict-aware rule (Eq.~6); conflict mass is
%     reassigned to ignorance, not renormalised. The pignistic transform
%     (Eq.~7) converts the fused mass to a scalar Decision Trust score
%     mapped to Accept / Caution / Reject (Eq.~8).}
%     \label{fig:evidence_fusion}
%   \end{figure*}
% ============================================================

\begin{tikzpicture}[
  every node/.append style={font=\sffamily},
  arrow/.style={-{Stealth[length=5pt]}, thick, gray!55},
  srcblock/.style={
    rectangle, rounded corners=2pt,
    minimum width=2.0cm, minimum height=0.7cm,
    font=\scriptsize\bfseries, align=center,
    draw, thick
  },
  % [LAYOUT] Increased massblock minimum width from 2.4cm to 2.8cm so
  %          multi-line math content wraps cleanly without manual \\ breaks
  %          Increased minimum height to 1.5cm for equation breathing room.
  massblock/.style={
    rectangle, rounded corners=2pt,
    minimum width=2.8cm, minimum height=1.5cm,
    font=\tiny, align=center,
    draw=gray!50, fill=gray!7, thin
  },
  % [LAYOUT] fuseblock sub-label font changed from embedded \tiny{} token
  %          (which caused font-command-in-alignment issues) to a separate
  %          node outside the main label. The fuseblock label itself uses
  %          the node's font=\scriptsize\bfseries consistently.
  fuseblock/.style={
    rectangle, rounded corners=3pt,
    minimum width=2.2cm, minimum height=1.2cm,
    font=\scriptsize\bfseries, align=center,
    draw=purple!60, fill=purple!8, thick
  },
  outblock/.style={
    rectangle, rounded corners=2pt,
    minimum width=1.5cm, minimum height=0.55cm,
    font=\scriptsize\bfseries, align=center,
    draw, thick
  },
  node distance=0.5cm and 0.55cm
]

\definecolor{pkilight}{RGB}{174,199,232}
\definecolor{pkidark}{RGB}{31,119,180}
\definecolor{mbdlight}{RGB}{152,223,138}
\definecolor{mbddark}{RGB}{44,160,44}
\definecolor{b3light}{RGB}{255,187,120}
\definecolor{b3dark}{RGB}{255,127,14}
\definecolor{fusepurple}{RGB}{148,103,189}
\definecolor{accgreen}{RGB}{44,160,44}
\definecolor{cauorange}{RGB}{255,127,14}
\definecolor{rejred}{RGB}{214,39,40}

% ============================================================
% ROW 1: Source blocks
% ============================================================
\node[srcblock, fill=pkilight, draw=pkidark, text=pkidark]
  (pki) {PKI / B1\\(crypto)};

\node[srcblock, fill=mbdlight, draw=mbddark, text=mbddark,
      right=1.0cm of pki]
  (mbd) {MBD / B2\\(behavioral)};

\node[srcblock, fill=mbdlight!60, draw=mbddark!60, text=mbddark,
      right=1.0cm of mbd]
  (cp) {CP\\(cooperative)};

\node[srcblock, fill=b3light, draw=b3dark, text=b3dark,
      right=1.0cm of cp]
  (b3) {B3\\(semantic)};

% ============================================================
% ROW 2: Mass function blocks
% [LAYOUT] Manual \\ inside \textit{} replaced with text-width
%          wrapping now that massblock is 2.8cm wide.
%          The math expressions still use \\ for alignment (correct
%          usage inside equation-style content).
% ============================================================
\node[massblock, below=0.6cm of pki] (mpki) {
  $m(T) = s_1 \cdot c_1$\\[4pt]
  $m(\bar{T}) = (1-s_1)c_1$\\[4pt]
  $m(\Theta) = 1-c_1$\\[6pt]
  \textit{e.g.} $s{=}1.0,\; c{=}0.46$
};

\node[massblock, below=0.6cm of mbd] (mmbd) {
  $m(T) = s_2 \cdot c_2$\\[4pt]
  $m(\bar{T}) = (1-s_2)c_2$\\[4pt]
  $m(\Theta) = 1-c_2$\\[6pt]
  \textit{behavioral score}
};

% [LAYOUT] Removed manual \\ from \textit{vacuous if no} line;
%          2.8cm massblock width allows it to wrap naturally.
\node[massblock, below=0.6cm of cp] (mcp) {
  $m(\Theta) = 1.0$\\[6pt]
  \textit{vacuous; no event label}\\[4pt]
  (no evidence)
};

\node[massblock, below=0.6cm of b3] (mb3) {
  $m(T) = s_3 \cdot c_3$\\[4pt]
  $m(\bar{T}) = (1-s_3)c_3$\\[4pt]
  $m(\Theta) = 1-c_3$\\[6pt]
  \textit{e.g.} $c{=}0.699$ (med.\ risk)
};

% ============================================================
% ARROWS: sources -> mass blocks
% ============================================================
\draw[arrow, pkidark!60]  (pki.south) -- (mpki.north);
\draw[arrow, mbddark!60]  (mbd.south) -- (mmbd.north);
\draw[arrow, mbddark!40]  (cp.south)  -- (mcp.north);
\draw[arrow, b3dark!60]   (b3.south)  -- (mb3.north);

% ============================================================
% ROW 3: Sequential fusion via Yager
% [LAYOUT] Fuseblock sub-labels now use \textit{\tiny ...} placed
%          as a separate line inside the node via \\ to avoid
%          embedded \tiny command issues. The label is kept as one
%          multi-line string within the node text.
% ============================================================

% First fusion: PKI + MBD
\node[fuseblock,
  below=1.2cm of mpki, xshift=0.5cm
] (fuse1) {Yager\\Fusion\\(Eq.~6)\\\textit{\fontsize{6}{7}\selectfont crypto+behavioral}};

\draw[arrow, fusepurple!60] (mpki.south) -- ([xshift=-0.5cm]fuse1.north);
\draw[arrow, fusepurple!60] (mmbd.south) |- ([xshift=0.5cm, yshift=0.3cm]fuse1.north) -- ([xshift=0.5cm]fuse1.north);

% Conflict annotation for fuse1
\node[font=\tiny, text=fusepurple!80, align=center,
      below=0.08cm of fuse1, yshift=-0.05cm]
  {conflict $K \to m(\Theta)$};

% Second fusion: result + CP (vacuous, no change)
\node[fuseblock,
  right=0.7cm of fuse1
] (fuse2) {Yager\\Fusion\\(Eq.~6)\\\textit{\fontsize{6}{7}\selectfont +~CP}};

\draw[arrow, fusepurple!60] (fuse1.east) -- (fuse2.west);
\draw[arrow, fusepurple!40] (mcp.south) |- ([yshift=0.6cm]fuse2.north) -- (fuse2.north);

% Third fusion: result + B3
\node[fuseblock,
  right=0.7cm of fuse2
] (fuse3) {Yager\\Fusion\\(Eq.~6)\\\textit{\fontsize{6}{7}\selectfont +~B3 semantic}};

\draw[arrow, fusepurple!60] (fuse2.east) -- (fuse3.west);
\draw[arrow, b3dark!60]     (mb3.south) |- ([yshift=0.9cm]fuse3.north) -- (fuse3.north);

% ============================================================
% ROW 4: Pignistic transform -> Decision Trust score
% ============================================================
\node[
  rectangle, rounded corners=3pt,
  below=0.7cm of fuse3,
  minimum width=2.6cm, minimum height=0.8cm,
  font=\scriptsize\bfseries, align=center,
  draw=gray!60, fill=gray!8, thick
] (pign) {Pignistic Transform\\(Eq.~7)\\$T_\text{Decision} = m_Y(T) + \tfrac{1}{2}m_Y(\Theta)$};

\draw[arrow] (fuse3.south) -- (pign.north);

% Also feed fuse2 result down (shown as continue arrow)
\draw[arrow, fusepurple!40] (fuse1.south) -- ++(0,-0.3)
  node[font=\tiny, text=fusepurple!70, anchor=west, xshift=0.05cm]
  {};

% ============================================================
% ROW 5: Three-way decision
% ============================================================
\node[outblock, fill=accgreen!20, draw=accgreen, text=accgreen,
      below left=0.6cm and -0.5cm of pign]
  (acc) {Accept\\$\geq 0.70$};

\node[outblock, fill=cauorange!20, draw=cauorange, text=cauorange,
      below=0.6cm of pign]
  (cau) {Caution\\$[0.40, 0.70)$};

\node[outblock, fill=rejred!15, draw=rejred, text=rejred,
      below right=0.6cm and -0.5cm of pign]
  (rej) {Reject\\$< 0.40$};

\draw[arrow, accgreen!60] (pign.south) -- ++(0,-0.18) -| (acc.north);
\draw[arrow, cauorange!60] (pign.south) -- (cau.north);
\draw[arrow, rejred!60] (pign.south) -- ++(0,-0.18) -| (rej.north);

% ============================================================
% FLOOR RULE ANNOTATION
% [LAYOUT] Reduced text width from 3.8cm to 3.0cm and moved anchor
%          slightly to prevent overflow beyond figure right boundary.
%          Removed manual \\ line-breaks; text width handles wrapping.
% ============================================================
\node[
  font=\tiny\itshape, text=b3dark, align=left, text width=4.5cm,
  right=0.3cm of rej
] {Floor rule: high-confidence B3 risk floors decision at Reject
   regardless of crypto/behavioral evidence.
   Semantic evidence only adds caution, never relaxes a rejection.};

% ============================================================
% WORKED EXAMPLE CALLOUT
% ============================================================
\node[
  rectangle, rounded corners=2pt,
  below=0.5cm of acc, xshift=0.5cm,
  minimum width=5.5cm, minimum height=0.7cm,
  font=\tiny, align=left,
  draw=gray!40, fill=gray!5, thin
] (example) {
  \textbf{Appendix B example:}\;
  $m_A{=}0.14,\; m_{\bar{A}}{=}0.38,\; m_\Theta{=}0.48$
  \;$\Rightarrow$\; $T_\text{Decision}{=}0.381$
  \;$\Rightarrow$\; \textcolor{rejred}{\textbf{REJECT}}
};

\end{tikzpicture}

%     \caption{Dempster--Shafer evidence fusion pipeline. Each evidence
%     source is converted to a basic belief assignment over
%     $\Theta=\{T,\bar{T}\}$ (Eq.~4). Sources are fused sequentially
%     via Yager's conflict-aware rule (Eq.~6); conflict mass is
%     reassigned to ignorance, not renormalised. The pignistic transform
%     (Eq.~7) converts the fused mass to a scalar Decision Trust score
%     mapped to Accept / Caution / Reject (Eq.~8).}
%     \label{fig:evidence_fusion}
%   \end{figure*}
% ============================================================

\begin{tikzpicture}[
  every node/.append style={font=\sffamily},
  arrow/.style={-{Stealth[length=5pt]}, thick, gray!55},
  srcblock/.style={
    rectangle, rounded corners=2pt,
    minimum width=2.0cm, minimum height=0.7cm,
    font=\scriptsize\bfseries, align=center,
    draw, thick
  },
  % [LAYOUT] Increased massblock minimum width from 2.4cm to 2.8cm so
  %          multi-line math content wraps cleanly without manual \\ breaks
  %          Increased minimum height to 1.5cm for equation breathing room.
  massblock/.style={
    rectangle, rounded corners=2pt,
    minimum width=2.8cm, minimum height=1.5cm,
    font=\tiny, align=center,
    draw=gray!50, fill=gray!7, thin
  },
  % [LAYOUT] fuseblock sub-label font changed from embedded \tiny{} token
  %          (which caused font-command-in-alignment issues) to a separate
  %          node outside the main label. The fuseblock label itself uses
  %          the node's font=\scriptsize\bfseries consistently.
  fuseblock/.style={
    rectangle, rounded corners=3pt,
    minimum width=2.2cm, minimum height=1.2cm,
    font=\scriptsize\bfseries, align=center,
    draw=purple!60, fill=purple!8, thick
  },
  outblock/.style={
    rectangle, rounded corners=2pt,
    minimum width=1.5cm, minimum height=0.55cm,
    font=\scriptsize\bfseries, align=center,
    draw, thick
  },
  node distance=0.5cm and 0.55cm
]

\definecolor{pkilight}{RGB}{174,199,232}
\definecolor{pkidark}{RGB}{31,119,180}
\definecolor{mbdlight}{RGB}{152,223,138}
\definecolor{mbddark}{RGB}{44,160,44}
\definecolor{b3light}{RGB}{255,187,120}
\definecolor{b3dark}{RGB}{255,127,14}
\definecolor{fusepurple}{RGB}{148,103,189}
\definecolor{accgreen}{RGB}{44,160,44}
\definecolor{cauorange}{RGB}{255,127,14}
\definecolor{rejred}{RGB}{214,39,40}

% ============================================================
% ROW 1: Source blocks
% ============================================================
\node[srcblock, fill=pkilight, draw=pkidark, text=pkidark]
  (pki) {PKI / B1\\(crypto)};

\node[srcblock, fill=mbdlight, draw=mbddark, text=mbddark,
      right=1.0cm of pki]
  (mbd) {MBD / B2\\(behavioral)};

\node[srcblock, fill=mbdlight!60, draw=mbddark!60, text=mbddark,
      right=1.0cm of mbd]
  (cp) {CP\\(cooperative)};

\node[srcblock, fill=b3light, draw=b3dark, text=b3dark,
      right=1.0cm of cp]
  (b3) {B3\\(semantic)};

% ============================================================
% ROW 2: Mass function blocks
% [LAYOUT] Manual \\ inside \textit{} replaced with text-width
%          wrapping now that massblock is 2.8cm wide.
%          The math expressions still use \\ for alignment (correct
%          usage inside equation-style content).
% ============================================================
\node[massblock, below=0.6cm of pki] (mpki) {
  $m(T) = s_1 \cdot c_1$\\[4pt]
  $m(\bar{T}) = (1-s_1)c_1$\\[4pt]
  $m(\Theta) = 1-c_1$\\[6pt]
  \textit{e.g.} $s{=}1.0,\; c{=}0.46$
};

\node[massblock, below=0.6cm of mbd] (mmbd) {
  $m(T) = s_2 \cdot c_2$\\[4pt]
  $m(\bar{T}) = (1-s_2)c_2$\\[4pt]
  $m(\Theta) = 1-c_2$\\[6pt]
  \textit{behavioral score}
};

% [LAYOUT] Removed manual \\ from \textit{vacuous if no} line;
%          2.8cm massblock width allows it to wrap naturally.
\node[massblock, below=0.6cm of cp] (mcp) {
  $m(\Theta) = 1.0$\\[6pt]
  \textit{vacuous; no event label}\\[4pt]
  (no evidence)
};

\node[massblock, below=0.6cm of b3] (mb3) {
  $m(T) = s_3 \cdot c_3$\\[4pt]
  $m(\bar{T}) = (1-s_3)c_3$\\[4pt]
  $m(\Theta) = 1-c_3$\\[6pt]
  \textit{e.g.} $c{=}0.699$ (med.\ risk)
};

% ============================================================
% ARROWS: sources -> mass blocks
% ============================================================
\draw[arrow, pkidark!60]  (pki.south) -- (mpki.north);
\draw[arrow, mbddark!60]  (mbd.south) -- (mmbd.north);
\draw[arrow, mbddark!40]  (cp.south)  -- (mcp.north);
\draw[arrow, b3dark!60]   (b3.south)  -- (mb3.north);

% ============================================================
% ROW 3: Sequential fusion via Yager
% [LAYOUT] Fuseblock sub-labels now use \textit{\tiny ...} placed
%          as a separate line inside the node via \\ to avoid
%          embedded \tiny command issues. The label is kept as one
%          multi-line string within the node text.
% ============================================================

% First fusion: PKI + MBD
\node[fuseblock,
  below=1.2cm of mpki, xshift=0.5cm
] (fuse1) {Yager\\Fusion\\(Eq.~6)\\\textit{\fontsize{6}{7}\selectfont crypto+behavioral}};

\draw[arrow, fusepurple!60] (mpki.south) -- ([xshift=-0.5cm]fuse1.north);
\draw[arrow, fusepurple!60] (mmbd.south) |- ([xshift=0.5cm, yshift=0.3cm]fuse1.north) -- ([xshift=0.5cm]fuse1.north);

% Conflict annotation for fuse1
\node[font=\tiny, text=fusepurple!80, align=center,
      below=0.08cm of fuse1, yshift=-0.05cm]
  {conflict $K \to m(\Theta)$};

% Second fusion: result + CP (vacuous, no change)
\node[fuseblock,
  right=0.7cm of fuse1
] (fuse2) {Yager\\Fusion\\(Eq.~6)\\\textit{\fontsize{6}{7}\selectfont +~CP}};

\draw[arrow, fusepurple!60] (fuse1.east) -- (fuse2.west);
\draw[arrow, fusepurple!40] (mcp.south) |- ([yshift=0.6cm]fuse2.north) -- (fuse2.north);

% Third fusion: result + B3
\node[fuseblock,
  right=0.7cm of fuse2
] (fuse3) {Yager\\Fusion\\(Eq.~6)\\\textit{\fontsize{6}{7}\selectfont +~B3 semantic}};

\draw[arrow, fusepurple!60] (fuse2.east) -- (fuse3.west);
\draw[arrow, b3dark!60]     (mb3.south) |- ([yshift=0.9cm]fuse3.north) -- (fuse3.north);

% ============================================================
% ROW 4: Pignistic transform -> Decision Trust score
% ============================================================
\node[
  rectangle, rounded corners=3pt,
  below=0.7cm of fuse3,
  minimum width=2.6cm, minimum height=0.8cm,
  font=\scriptsize\bfseries, align=center,
  draw=gray!60, fill=gray!8, thick
] (pign) {Pignistic Transform\\(Eq.~7)\\$T_\text{Decision} = m_Y(T) + \tfrac{1}{2}m_Y(\Theta)$};

\draw[arrow] (fuse3.south) -- (pign.north);

% Also feed fuse2 result down (shown as continue arrow)
\draw[arrow, fusepurple!40] (fuse1.south) -- ++(0,-0.3)
  node[font=\tiny, text=fusepurple!70, anchor=west, xshift=0.05cm]
  {};

% ============================================================
% ROW 5: Three-way decision
% ============================================================
\node[outblock, fill=accgreen!20, draw=accgreen, text=accgreen,
      below left=0.6cm and -0.5cm of pign]
  (acc) {Accept\\$\geq 0.70$};

\node[outblock, fill=cauorange!20, draw=cauorange, text=cauorange,
      below=0.6cm of pign]
  (cau) {Caution\\$[0.40, 0.70)$};

\node[outblock, fill=rejred!15, draw=rejred, text=rejred,
      below right=0.6cm and -0.5cm of pign]
  (rej) {Reject\\$< 0.40$};

\draw[arrow, accgreen!60] (pign.south) -- ++(0,-0.18) -| (acc.north);
\draw[arrow, cauorange!60] (pign.south) -- (cau.north);
\draw[arrow, rejred!60] (pign.south) -- ++(0,-0.18) -| (rej.north);

% ============================================================
% FLOOR RULE ANNOTATION
% [LAYOUT] Reduced text width from 3.8cm to 3.0cm and moved anchor
%          slightly to prevent overflow beyond figure right boundary.
%          Removed manual \\ line-breaks; text width handles wrapping.
% ============================================================
\node[
  font=\tiny\itshape, text=b3dark, align=left, text width=4.5cm,
  right=0.3cm of rej
] {Floor rule: high-confidence B3 risk floors decision at Reject
   regardless of crypto/behavioral evidence.
   Semantic evidence only adds caution, never relaxes a rejection.};

% ============================================================
% WORKED EXAMPLE CALLOUT
% ============================================================
\node[
  rectangle, rounded corners=2pt,
  below=0.5cm of acc, xshift=0.5cm,
  minimum width=5.5cm, minimum height=0.7cm,
  font=\tiny, align=left,
  draw=gray!40, fill=gray!5, thin
] (example) {
  \textbf{Appendix B example:}\;
  $m_A{=}0.14,\; m_{\bar{A}}{=}0.38,\; m_\Theta{=}0.48$
  \;$\Rightarrow$\; $T_\text{Decision}{=}0.381$
  \;$\Rightarrow$\; \textcolor{rejred}{\textbf{REJECT}}
};

\end{tikzpicture}

%     \caption{Dempster--Shafer evidence fusion pipeline. Each evidence
%     source is converted to a basic belief assignment over
%     $\Theta=\{T,\bar{T}\}$ (Eq.~4). Sources are fused sequentially
%     via Yager's conflict-aware rule (Eq.~6); conflict mass is
%     reassigned to ignorance, not renormalised. The pignistic transform
%     (Eq.~7) converts the fused mass to a scalar Decision Trust score
%     mapped to Accept / Caution / Reject (Eq.~8).}
%     \label{fig:evidence_fusion}
%   \end{figure*}
% ============================================================

\begin{tikzpicture}[
  every node/.append style={font=\sffamily},
  arrow/.style={-{Stealth[length=5pt]}, thick, gray!55},
  srcblock/.style={
    rectangle, rounded corners=2pt,
    minimum width=2.0cm, minimum height=0.7cm,
    font=\scriptsize\bfseries, align=center,
    draw, thick
  },
  % [LAYOUT] Increased massblock minimum width from 2.4cm to 2.8cm so
  %          multi-line math content wraps cleanly without manual \\ breaks
  %          Increased minimum height to 1.5cm for equation breathing room.
  massblock/.style={
    rectangle, rounded corners=2pt,
    minimum width=2.8cm, minimum height=1.5cm,
    font=\tiny, align=center,
    draw=gray!50, fill=gray!7, thin
  },
  % [LAYOUT] fuseblock sub-label font changed from embedded \tiny{} token
  %          (which caused font-command-in-alignment issues) to a separate
  %          node outside the main label. The fuseblock label itself uses
  %          the node's font=\scriptsize\bfseries consistently.
  fuseblock/.style={
    rectangle, rounded corners=3pt,
    minimum width=2.2cm, minimum height=1.2cm,
    font=\scriptsize\bfseries, align=center,
    draw=purple!60, fill=purple!8, thick
  },
  outblock/.style={
    rectangle, rounded corners=2pt,
    minimum width=1.5cm, minimum height=0.55cm,
    font=\scriptsize\bfseries, align=center,
    draw, thick
  },
  node distance=0.5cm and 0.55cm
]

\definecolor{pkilight}{RGB}{174,199,232}
\definecolor{pkidark}{RGB}{31,119,180}
\definecolor{mbdlight}{RGB}{152,223,138}
\definecolor{mbddark}{RGB}{44,160,44}
\definecolor{b3light}{RGB}{255,187,120}
\definecolor{b3dark}{RGB}{255,127,14}
\definecolor{fusepurple}{RGB}{148,103,189}
\definecolor{accgreen}{RGB}{44,160,44}
\definecolor{cauorange}{RGB}{255,127,14}
\definecolor{rejred}{RGB}{214,39,40}

% ============================================================
% ROW 1: Source blocks
% ============================================================
\node[srcblock, fill=pkilight, draw=pkidark, text=pkidark]
  (pki) {PKI / B1\\(crypto)};

\node[srcblock, fill=mbdlight, draw=mbddark, text=mbddark,
      right=1.0cm of pki]
  (mbd) {MBD / B2\\(behavioral)};

\node[srcblock, fill=mbdlight!60, draw=mbddark!60, text=mbddark,
      right=1.0cm of mbd]
  (cp) {CP\\(cooperative)};

\node[srcblock, fill=b3light, draw=b3dark, text=b3dark,
      right=1.0cm of cp]
  (b3) {B3\\(semantic)};

% ============================================================
% ROW 2: Mass function blocks
% [LAYOUT] Manual \\ inside \textit{} replaced with text-width
%          wrapping now that massblock is 2.8cm wide.
%          The math expressions still use \\ for alignment (correct
%          usage inside equation-style content).
% ============================================================
\node[massblock, below=0.6cm of pki] (mpki) {
  $m(T) = s_1 \cdot c_1$\\[4pt]
  $m(\bar{T}) = (1-s_1)c_1$\\[4pt]
  $m(\Theta) = 1-c_1$\\[6pt]
  \textit{e.g.} $s{=}1.0,\; c{=}0.46$
};

\node[massblock, below=0.6cm of mbd] (mmbd) {
  $m(T) = s_2 \cdot c_2$\\[4pt]
  $m(\bar{T}) = (1-s_2)c_2$\\[4pt]
  $m(\Theta) = 1-c_2$\\[6pt]
  \textit{behavioral score}
};

% [LAYOUT] Removed manual \\ from \textit{vacuous if no} line;
%          2.8cm massblock width allows it to wrap naturally.
\node[massblock, below=0.6cm of cp] (mcp) {
  $m(\Theta) = 1.0$\\[6pt]
  \textit{vacuous; no event label}\\[4pt]
  (no evidence)
};

\node[massblock, below=0.6cm of b3] (mb3) {
  $m(T) = s_3 \cdot c_3$\\[4pt]
  $m(\bar{T}) = (1-s_3)c_3$\\[4pt]
  $m(\Theta) = 1-c_3$\\[6pt]
  \textit{e.g.} $c{=}0.699$ (med.\ risk)
};

% ============================================================
% ARROWS: sources -> mass blocks
% ============================================================
\draw[arrow, pkidark!60]  (pki.south) -- (mpki.north);
\draw[arrow, mbddark!60]  (mbd.south) -- (mmbd.north);
\draw[arrow, mbddark!40]  (cp.south)  -- (mcp.north);
\draw[arrow, b3dark!60]   (b3.south)  -- (mb3.north);

% ============================================================
% ROW 3: Sequential fusion via Yager
% [LAYOUT] Fuseblock sub-labels now use \textit{\tiny ...} placed
%          as a separate line inside the node via \\ to avoid
%          embedded \tiny command issues. The label is kept as one
%          multi-line string within the node text.
% ============================================================

% First fusion: PKI + MBD
\node[fuseblock,
  below=1.2cm of mpki, xshift=0.5cm
] (fuse1) {Yager\\Fusion\\(Eq.~6)\\\textit{\fontsize{6}{7}\selectfont crypto+behavioral}};

\draw[arrow, fusepurple!60] (mpki.south) -- ([xshift=-0.5cm]fuse1.north);
\draw[arrow, fusepurple!60] (mmbd.south) |- ([xshift=0.5cm, yshift=0.3cm]fuse1.north) -- ([xshift=0.5cm]fuse1.north);

% Conflict annotation for fuse1
\node[font=\tiny, text=fusepurple!80, align=center,
      below=0.08cm of fuse1, yshift=-0.05cm]
  {conflict $K \to m(\Theta)$};

% Second fusion: result + CP (vacuous, no change)
\node[fuseblock,
  right=0.7cm of fuse1
] (fuse2) {Yager\\Fusion\\(Eq.~6)\\\textit{\fontsize{6}{7}\selectfont +~CP}};

\draw[arrow, fusepurple!60] (fuse1.east) -- (fuse2.west);
\draw[arrow, fusepurple!40] (mcp.south) |- ([yshift=0.6cm]fuse2.north) -- (fuse2.north);

% Third fusion: result + B3
\node[fuseblock,
  right=0.7cm of fuse2
] (fuse3) {Yager\\Fusion\\(Eq.~6)\\\textit{\fontsize{6}{7}\selectfont +~B3 semantic}};

\draw[arrow, fusepurple!60] (fuse2.east) -- (fuse3.west);
\draw[arrow, b3dark!60]     (mb3.south) |- ([yshift=0.9cm]fuse3.north) -- (fuse3.north);

% ============================================================
% ROW 4: Pignistic transform -> Decision Trust score
% ============================================================
\node[
  rectangle, rounded corners=3pt,
  below=0.7cm of fuse3,
  minimum width=2.6cm, minimum height=0.8cm,
  font=\scriptsize\bfseries, align=center,
  draw=gray!60, fill=gray!8, thick
] (pign) {Pignistic Transform\\(Eq.~7)\\$T_\text{Decision} = m_Y(T) + \tfrac{1}{2}m_Y(\Theta)$};

\draw[arrow] (fuse3.south) -- (pign.north);

% Also feed fuse2 result down (shown as continue arrow)
\draw[arrow, fusepurple!40] (fuse1.south) -- ++(0,-0.3)
  node[font=\tiny, text=fusepurple!70, anchor=west, xshift=0.05cm]
  {};

% ============================================================
% ROW 5: Three-way decision
% ============================================================
\node[outblock, fill=accgreen!20, draw=accgreen, text=accgreen,
      below left=0.6cm and -0.5cm of pign]
  (acc) {Accept\\$\geq 0.70$};

\node[outblock, fill=cauorange!20, draw=cauorange, text=cauorange,
      below=0.6cm of pign]
  (cau) {Caution\\$[0.40, 0.70)$};

\node[outblock, fill=rejred!15, draw=rejred, text=rejred,
      below right=0.6cm and -0.5cm of pign]
  (rej) {Reject\\$< 0.40$};

\draw[arrow, accgreen!60] (pign.south) -- ++(0,-0.18) -| (acc.north);
\draw[arrow, cauorange!60] (pign.south) -- (cau.north);
\draw[arrow, rejred!60] (pign.south) -- ++(0,-0.18) -| (rej.north);

% ============================================================
% FLOOR RULE ANNOTATION
% [LAYOUT] Reduced text width from 3.8cm to 3.0cm and moved anchor
%          slightly to prevent overflow beyond figure right boundary.
%          Removed manual \\ line-breaks; text width handles wrapping.
% ============================================================
\node[
  font=\tiny\itshape, text=b3dark, align=left, text width=4.5cm,
  right=0.3cm of rej
] {Floor rule: high-confidence B3 risk floors decision at Reject
   regardless of crypto/behavioral evidence.
   Semantic evidence only adds caution, never relaxes a rejection.};

% ============================================================
% WORKED EXAMPLE CALLOUT
% ============================================================
\node[
  rectangle, rounded corners=2pt,
  below=0.5cm of acc, xshift=0.5cm,
  minimum width=5.5cm, minimum height=0.7cm,
  font=\tiny, align=left,
  draw=gray!40, fill=gray!5, thin
] (example) {
  \textbf{Appendix B example:}\;
  $m_A{=}0.14,\; m_{\bar{A}}{=}0.38,\; m_\Theta{=}0.48$
  \;$\Rightarrow$\; $T_\text{Decision}{=}0.381$
  \;$\Rightarrow$\; \textcolor{rejred}{\textbf{REJECT}}
};

\end{tikzpicture}

%     \caption{Dempster--Shafer evidence fusion pipeline. Each evidence
%     source is converted to a basic belief assignment over
%     $\Theta=\{T,\bar{T}\}$ (Eq.~4). Sources are fused sequentially
%     via Yager's conflict-aware rule (Eq.~6); conflict mass is
%     reassigned to ignorance, not renormalised. The pignistic transform
%     (Eq.~7) converts the fused mass to a scalar Decision Trust score
%     mapped to Accept / Caution / Reject (Eq.~8).}
%     \label{fig:evidence_fusion}
%   \end{figure*}
% ============================================================

\begin{tikzpicture}[
  every node/.append style={font=\sffamily},
  arrow/.style={-{Stealth[length=5pt]}, thick, gray!55},
  srcblock/.style={
    rectangle, rounded corners=2pt,
    minimum width=2.0cm, minimum height=0.7cm,
    font=\scriptsize\bfseries, align=center,
    draw, thick
  },
  % [LAYOUT] Increased massblock minimum width from 2.4cm to 2.8cm so
  %          multi-line math content wraps cleanly without manual \\ breaks
  %          Increased minimum height to 1.5cm for equation breathing room.
  massblock/.style={
    rectangle, rounded corners=2pt,
    minimum width=2.8cm, minimum height=1.5cm,
    font=\tiny, align=center,
    draw=gray!50, fill=gray!7, thin
  },
  % [LAYOUT] fuseblock sub-label font changed from embedded \tiny{} token
  %          (which caused font-command-in-alignment issues) to a separate
  %          node outside the main label. The fuseblock label itself uses
  %          the node's font=\scriptsize\bfseries consistently.
  fuseblock/.style={
    rectangle, rounded corners=3pt,
    minimum width=2.2cm, minimum height=1.2cm,
    font=\scriptsize\bfseries, align=center,
    draw=purple!60, fill=purple!8, thick
  },
  outblock/.style={
    rectangle, rounded corners=2pt,
    minimum width=1.5cm, minimum height=0.55cm,
    font=\scriptsize\bfseries, align=center,
    draw, thick
  },
  node distance=0.5cm and 0.55cm
]

\definecolor{pkilight}{RGB}{174,199,232}
\definecolor{pkidark}{RGB}{31,119,180}
\definecolor{mbdlight}{RGB}{152,223,138}
\definecolor{mbddark}{RGB}{44,160,44}
\definecolor{b3light}{RGB}{255,187,120}
\definecolor{b3dark}{RGB}{255,127,14}
\definecolor{fusepurple}{RGB}{148,103,189}
\definecolor{accgreen}{RGB}{44,160,44}
\definecolor{cauorange}{RGB}{255,127,14}
\definecolor{rejred}{RGB}{214,39,40}

% ============================================================
% ROW 1: Source blocks
% ============================================================
\node[srcblock, fill=pkilight, draw=pkidark, text=pkidark]
  (pki) {PKI / B1\\(crypto)};

\node[srcblock, fill=mbdlight, draw=mbddark, text=mbddark,
      right=1.0cm of pki]
  (mbd) {MBD / B2\\(behavioral)};

\node[srcblock, fill=mbdlight!60, draw=mbddark!60, text=mbddark,
      right=1.0cm of mbd]
  (cp) {CP\\(cooperative)};

\node[srcblock, fill=b3light, draw=b3dark, text=b3dark,
      right=1.0cm of cp]
  (b3) {B3\\(semantic)};

% ============================================================
% ROW 2: Mass function blocks
% [LAYOUT] Manual \\ inside \textit{} replaced with text-width
%          wrapping now that massblock is 2.8cm wide.
%          The math expressions still use \\ for alignment (correct
%          usage inside equation-style content).
% ============================================================
\node[massblock, below=0.6cm of pki] (mpki) {
  $m(T) = s_1 \cdot c_1$\\[4pt]
  $m(\bar{T}) = (1-s_1)c_1$\\[4pt]
  $m(\Theta) = 1-c_1$\\[6pt]
  \textit{e.g.} $s{=}1.0,\; c{=}0.46$
};

\node[massblock, below=0.6cm of mbd] (mmbd) {
  $m(T) = s_2 \cdot c_2$\\[4pt]
  $m(\bar{T}) = (1-s_2)c_2$\\[4pt]
  $m(\Theta) = 1-c_2$\\[6pt]
  \textit{behavioral score}
};

% [LAYOUT] Removed manual \\ from \textit{vacuous if no} line;
%          2.8cm massblock width allows it to wrap naturally.
\node[massblock, below=0.6cm of cp] (mcp) {
  $m(\Theta) = 1.0$\\[6pt]
  \textit{vacuous; no event label}\\[4pt]
  (no evidence)
};

\node[massblock, below=0.6cm of b3] (mb3) {
  $m(T) = s_3 \cdot c_3$\\[4pt]
  $m(\bar{T}) = (1-s_3)c_3$\\[4pt]
  $m(\Theta) = 1-c_3$\\[6pt]
  \textit{e.g.} $c{=}0.699$ (med.\ risk)
};

% ============================================================
% ARROWS: sources -> mass blocks
% ============================================================
\draw[arrow, pkidark!60]  (pki.south) -- (mpki.north);
\draw[arrow, mbddark!60]  (mbd.south) -- (mmbd.north);
\draw[arrow, mbddark!40]  (cp.south)  -- (mcp.north);
\draw[arrow, b3dark!60]   (b3.south)  -- (mb3.north);

% ============================================================
% ROW 3: Sequential fusion via Yager
% [LAYOUT] Fuseblock sub-labels now use \textit{\tiny ...} placed
%          as a separate line inside the node via \\ to avoid
%          embedded \tiny command issues. The label is kept as one
%          multi-line string within the node text.
% ============================================================

% First fusion: PKI + MBD
\node[fuseblock,
  below=1.2cm of mpki, xshift=0.5cm
] (fuse1) {Yager\\Fusion\\(Eq.~6)\\\textit{\fontsize{6}{7}\selectfont crypto+behavioral}};

\draw[arrow, fusepurple!60] (mpki.south) -- ([xshift=-0.5cm]fuse1.north);
\draw[arrow, fusepurple!60] (mmbd.south) |- ([xshift=0.5cm, yshift=0.3cm]fuse1.north) -- ([xshift=0.5cm]fuse1.north);

% Conflict annotation for fuse1
\node[font=\tiny, text=fusepurple!80, align=center,
      below=0.08cm of fuse1, yshift=-0.05cm]
  {conflict $K \to m(\Theta)$};

% Second fusion: result + CP (vacuous, no change)
\node[fuseblock,
  right=0.7cm of fuse1
] (fuse2) {Yager\\Fusion\\(Eq.~6)\\\textit{\fontsize{6}{7}\selectfont +~CP}};

\draw[arrow, fusepurple!60] (fuse1.east) -- (fuse2.west);
\draw[arrow, fusepurple!40] (mcp.south) |- ([yshift=0.6cm]fuse2.north) -- (fuse2.north);

% Third fusion: result + B3
\node[fuseblock,
  right=0.7cm of fuse2
] (fuse3) {Yager\\Fusion\\(Eq.~6)\\\textit{\fontsize{6}{7}\selectfont +~B3 semantic}};

\draw[arrow, fusepurple!60] (fuse2.east) -- (fuse3.west);
\draw[arrow, b3dark!60]     (mb3.south) |- ([yshift=0.9cm]fuse3.north) -- (fuse3.north);

% ============================================================
% ROW 4: Pignistic transform -> Decision Trust score
% ============================================================
\node[
  rectangle, rounded corners=3pt,
  below=0.7cm of fuse3,
  minimum width=2.6cm, minimum height=0.8cm,
  font=\scriptsize\bfseries, align=center,
  draw=gray!60, fill=gray!8, thick
] (pign) {Pignistic Transform\\(Eq.~7)\\$T_\text{Decision} = m_Y(T) + \tfrac{1}{2}m_Y(\Theta)$};

\draw[arrow] (fuse3.south) -- (pign.north);

% Also feed fuse2 result down (shown as continue arrow)
\draw[arrow, fusepurple!40] (fuse1.south) -- ++(0,-0.3)
  node[font=\tiny, text=fusepurple!70, anchor=west, xshift=0.05cm]
  {};

% ============================================================
% ROW 5: Three-way decision
% ============================================================
\node[outblock, fill=accgreen!20, draw=accgreen, text=accgreen,
      below left=0.6cm and -0.5cm of pign]
  (acc) {Accept\\$\geq 0.70$};

\node[outblock, fill=cauorange!20, draw=cauorange, text=cauorange,
      below=0.6cm of pign]
  (cau) {Caution\\$[0.40, 0.70)$};

\node[outblock, fill=rejred!15, draw=rejred, text=rejred,
      below right=0.6cm and -0.5cm of pign]
  (rej) {Reject\\$< 0.40$};

\draw[arrow, accgreen!60] (pign.south) -- ++(0,-0.18) -| (acc.north);
\draw[arrow, cauorange!60] (pign.south) -- (cau.north);
\draw[arrow, rejred!60] (pign.south) -- ++(0,-0.18) -| (rej.north);

% ============================================================
% FLOOR RULE ANNOTATION
% [LAYOUT] Reduced text width from 3.8cm to 3.0cm and moved anchor
%          slightly to prevent overflow beyond figure right boundary.
%          Removed manual \\ line-breaks; text width handles wrapping.
% ============================================================
\node[
  font=\tiny\itshape, text=b3dark, align=left, text width=4.5cm,
  right=0.3cm of rej
] {Floor rule: high-confidence B3 risk floors decision at Reject
   regardless of crypto/behavioral evidence.
   Semantic evidence only adds caution, never relaxes a rejection.};

% ============================================================
% WORKED EXAMPLE CALLOUT
% ============================================================
\node[
  rectangle, rounded corners=2pt,
  below=0.5cm of acc, xshift=0.5cm,
  minimum width=5.5cm, minimum height=0.7cm,
  font=\tiny, align=left,
  draw=gray!40, fill=gray!5, thin
] (example) {
  \textbf{Appendix B example:}\;
  $m_A{=}0.14,\; m_{\bar{A}}{=}0.38,\; m_\Theta{=}0.48$
  \;$\Rightarrow$\; $T_\text{Decision}{=}0.381$
  \;$\Rightarrow$\; \textcolor{rejred}{\textbf{REJECT}}
};

\end{tikzpicture}
